% Source: Weinberg GR, physical PDF p. 395 (printed p. 375).
% Coverage: Chapter 13 opening material before Section 13.1.
% Figures/tables/footnotes: chapter epigraph; no figures, tables, or footnotes.
% Status: source-reviewed and compile-clean.
% Uncertainties: none.

\begin{flushright}
\begin{minipage}{0.46\textwidth}
\itshape
``Symmetry, as wide or as\\
narrow as you may define\\
it, is one idea by which man\\
through the ages has tried\\
to comprehend and create\\
order, beauty, and\\
perfection.'' Hermann Weyl,\\
\emph{Symmetry}
\end{minipage}
\end{flushright}

Euclid implicitly assumed that metric relations are unaffected by translations
or rotations. Real gravitational fields do not usually have such a high degree
of symmetry, but they often admit some group of approximate symmetry
transformations, and when they do, we can use this information to help solve
the Einstein equations, or even to do without a solution. I shall give only a
very brief introduction to the elaborate mathematical theory of symmetric
spaces, with special attention to the maximally symmetric spaces that are of
special interest in cosmology.

The initial difficulty here is: How can we use some supposed symmetry of a
metric space to gain information about the metric, when we need to know the
metric before we can establish a coordinate system in which to define the
symmetry? In order to avoid this impasse, we shall have to learn ways of
describing symmetries in a covariant language, which does not depend on any
particular choice of coordinate system. Once this language is established, it
becomes a matter of mathematical manipulation to determine those properties of
a metric that follow from its symmetries.
